\section{Related Work}

Emotion analysis is closely related to sentiment analysis, but it models a more detailed affective space. Sentiment analysis traditionally focuses on polarity, such as positive versus negative opinion \citep{pang2008opinion}. Emotion analysis instead asks which emotions are expressed or evoked by text. This distinction is important because two texts can both be negative while expressing different emotions, for example anger versus sadness or fear.

Lexicon-based emotion analysis has a long history in natural language processing. Emotion lexicons represent associations between words or phrases and affective categories. \citet{mohammad2010emotions} describe the creation of an emotion lexicon using crowdsourcing, and \citet{mohammad2013crowdsourcing} expand this work into the NRC Word--Emotion Association Lexicon. These resources are important because emotion-labeled text datasets can be expensive to create, while lexicons provide reusable affective knowledge.

The NRC Hashtag Emotion Lexicon is especially relevant for social-media-style text. It was automatically generated from tweets containing emotion-word hashtags, and provides weighted associations between terms and eight emotions \citep{mohammad2012emotional,mohammad2015hashtags}. This makes it suitable for a project that aims to detect emotions in short texts and social-media-like inputs. The lexicon contains rows of the form \texttt{emotion, term, score}, where the score indicates the strength of association between a term and an emotion.

Machine learning approaches provide another way to model emotion. Classic supervised algorithms, such as Naive Bayes, logistic regression, or support vector machines, can be trained on labeled examples. In this project, because a large manually annotated dataset was not used, the Naive Bayes model is weakly supervised: lexicon entries are converted into training examples, and lexical association scores become training weights. This does not replace a fully supervised neural model, but it provides a learned baseline that can be compared with the classic lexicon method.

Recent work in emotion analysis often uses neural models and Transformer architectures, especially when large labeled datasets are available. Such models can capture context better than isolated lexical lookup. However, they require more dependencies, often benefit from GPU acceleration, and are less transparent. The current implementation intentionally prioritizes reproducibility, interpretability, and local execution. A Transformer baseline is therefore treated as future work rather than a dependency of the submitted implementation.
